% ===== 问题三正文 =====
% 本文件汇总第三问；各部分正文、公式与图表请编辑 sections/q3/ 下的对应子文件。
% 算法依据 paper_code/q3_optical_paper 中的 v3 提前光学策略；实验数字只引用该包 code/statistics.json 的记录。
\subsection{问题三：多源自动搜索与清除策略}\label{sec:q3}

问题三要机器狗从原点出发，把场地里 $10$ 到 $16$ 个全向干扰源全部找到并清除，源的数量、位置和接收半径都不知道，目标是总耗时尽量短。和前两问相比，这里多了三件事要决定：还没发现的源怎么保证一定能找到；已经发现的源什么时候去定位、什么时候去清除；任务什么时候可以结束。这三件事互相牵制：多扫几站能早点把源找全，但每扫一站要花几十秒检测时间；先把手头的源清掉再去扫站，路可能要走回头。本节先交代沿用的记号和无信号反馈的用法，再分发现、定位、清除、任务安排四部分建立模型，其中定位部分在第二问的选点模型之外加了一条“区域已经很小时先试一次光学清除”的规则；然后给出求解流程；最后在本地复原的模拟器上用 $2000$ 个相同场景比较五种方案，检验各条规则的作用。

\subsubsection{模型准备}\label{sec:q3-1}

\textbf{（1）沿用的记号与信息状态。}目标圆域 $\Omega$、频道 $f$、源位置 $p$、接收半径 $R_f$、机器狗当前位置 $x_t$、定位区域的外包多边形 $\widehat C_f$ 及其包围圆的圆心 $z(\widehat C_f)$、半径 $\rho(\widehat C_f)$ 都沿用第二问的记号。程序把 $20$ 个频道分成四类：还没收到过信号的\textbf{未发现}频道；收到过示向度、正在定位的\textbf{定位中}频道，每个保存一个外包多边形、它的包围圆、测过的位置，以及收到过信号和没收到信号的位置 $x^+$、$x^-$；收到真实成功反馈的\textbf{已清除}频道；以及已经能断定没有源的\textbf{无源}频道。任务开始时 $20$ 个频道全在未发现一类，任务结束时全部落在已清除或无源两类。源的真实位置、数量和接收半径都不进入决策，算法只通过检测和清除的反馈逐步了解它们。

\textbf{（2）总耗时与评价指标。}按题目规则，一局的总耗时由五项累加：
\begin{equation}\label{eq:q3-2}
T=\frac{L}{v}+5M+C+3A+2K,\qquad v=5\ \text{m/s},
\end{equation}
其中 $L$ 是总路程，$M$ 是检测次数，$C$ 是换频次数，$A$ 是光学定位的尝试次数（成功和失败都算），$K$ 是成功清除的次数。所以一次成功的光学清除共花 $5$ s，一次失败只花 $3$ s；清除动作不改变测向机当前的频道。题目的评价指标是平均定位清除时间 $T/N$，$N$ 为清除的源数；后面比较不同方案时，先对每个场景算 $T/N$，再对场景等权平均。

\textbf{（3）无信号反馈的用法。}第二问只用到有信号的反馈：示向度给出一个 $\pm1^\circ$ 的楔形，与原区域求交。第三问里机器狗在每个检测站都要把未发现频道扫一遍，对已发现的源也常常顺便再测一次，因此会积累不少“在某处对某频道没有收到信号”的记录。对全向源，无信号只有一个原因：离源超过了接收半径。这一点能从两方面缩小定位区域，见图\ref{fig:q3-negative}。

第一，接收半径至少 $1000$ m，所以以无信号点 $x^-$ 为圆心、半径 $1000$ m 的圆盘内不可能有源。把这个圆盘从定位区域里挖掉以后，剩下的部分一般不再是凸的，程序取它的凸包作为新的外包多边形，这样只会放大、不会缩小区域，真实源不会被裁掉。

第二，同一个源的接收半径在一局里固定不变。若在 $x^+$ 收到了信号、在 $x^-$ 没有收到，则源到 $x^+$ 的距离不超过 $R_f$，到 $x^-$ 的距离超过 $R_f$，于是源离 $x^+$ 比离 $x^-$ 更近：
\begin{equation}\label{eq:q3-3}
\|p-x^+\|\le R_f<\|p-x^-\|
\;\Longrightarrow\;
(x^--x^+)^{\mathsf T}p\le\frac{\|x^-\|^2-\|x^+\|^2}{2}.
\end{equation}
右边是一个半平面，边界正是 $x^+x^-$ 的垂直平分线。每来一个新的无信号点，就把它和该频道所有有信号点配对，逐个切掉半平面；程序在右端加了 $10^{-5}$ 的余量，把严格不等式放宽成闭半平面。这两条推断都只对全向源成立，第四问的定向源不能这样用。

\begin{center}
\begin{minipage}{\linewidth}
\centering
\QThreeNegativeFigure
\captionof{figure}{无信号反馈的两种裁剪。(a) 以无信号点为圆心的 $1000$ m 圆盘内不可能有源，挖掉后取凸包。(b) 源到有信号点的距离不超过接收半径、到无信号点的距离超过接收半径，所以源在两点垂直平分线靠近 $x^+$ 的一侧。}
\label{fig:q3-negative}
\end{minipage}
\end{center}

\textbf{（4）定位区域的构造与更新。}某个频道第一次在位置 $x$ 收到示向度 $\theta$ 时，定位区域取三者的交：以 $x$ 为顶点、沿 $\theta$ 方向张开 $\pm\varepsilon$、长 $1500$ m 的三角形（它包含半径 $1500$ m 的扇形），场地的外切正 $64$ 边形，以及该频道此前所有无信号点按（3）给出的裁剪。以后每收到一次示向度，就再交上一个楔形，并重新做一遍无信号裁剪。程序里 $\varepsilon$ 取 $1.005^\circ$，比题面的 $1^\circ$ 多出的 $0.005^\circ$ 用来容纳示向度保留两位小数带来的舍入。每次更新后重新计算包围圆，方法是枚举多边形顶点中由一个、两个或三个顶点支撑的圆，取半径最小且包住全部顶点的那个；区域顶点很少，枚举很快。

\textbf{（5）不能省的三条底线。}后面所有的规则都在这三条之内做取舍：第一，真实源始终在它的定位区域里，区域只按真实反馈和上面两条推断收缩；第二，一个未发现频道只有在“它在场地里任何位置都会被某个已扫描的站收到”的条件下，才能判为无源；第三，只有收到设备返回的清除成功，才把频道记为已清除。程序把这三条写成运行时检查，任何一条被违反就报错停机，而不是继续跑下去。
%
\subsubsection{模型的建立}\label{sec:q3-2}

\textbf{（1）决策内容与目标。}要决定的是机器狗每一步去哪里、在那里测哪些频道、什么时候对哪个源做光学清除，目标是在清除全部源的前提下让式\eqref{eq:q3-2} 的总耗时尽量小。真实源的数量和位置事先不知道，每一次检测的结果又会改变后面该做什么，所以不可能一次把整局的路线规划出来。本文的做法是分层：先用一组固定的检测站保证任何源都会被发现；对已发现的源用定位区域描述它可能在哪里，按第二问的模型选点测向，区域很小时先试一次光学清除；最后把“去哪个站”和“处理哪个源”合成一个路线问题，每走一步重新规划。下面依次建立这四部分。

\textbf{（2）发现覆盖：$7$ 个检测站。}接收半径至少 $1000$ m，所以只要某个位置离站不到 $1000$ m，站在那里扫一遍所有未发现频道，这个位置上的源就一定会被收到。要保证场地里任何位置的源都能被发现，只需一组站，使它们的 $1000$ m 圆盘把整个半径 $1800$ m 的圆域盖住。本文取半径 $999$ m 的圆周上等角分布的 $7$ 个站，见图\ref{fig:q3-cover}(a)。相邻两站相隔 $360^\circ/7\approx51.4^\circ$，圆周上离站最远的点在两站中间，它到最近站的距离约 $999$ m，刚好在 $1000$ m 以内；原点离每个站都是 $999$ m，也在圆内；场地内部的点更近。程序不靠抽样检查覆盖，而是用圆弧关系做一次完整的核验：先检查场地圆周是否全被各站圆盘盖住，再对每个站检查它的圆在场地内的那段弧是否全被其他站的圆盘盖住。如果有一块地方没被盖住，它的边界上一定有某个站圆的一段露出来的弧，所以两条都通过就说明覆盖完整。

这个布局有一点和直觉不同：机器狗从原点出发，却不先在原点把 $20$ 个频道扫一遍。原因是原点在每个站的 $1000$ m 圆内，$7$ 个站扫过之后原点附近的源自然会被发现，而在原点多扫一遍要花 $20\times5+19$ s 左右。开局先扫原点的做法作为对照方案保留在结果分析里，实验表明它平均更慢。

固定的 $7$ 站是计划，不是必须逐一执行的。扫过几站以后，剩下的站可能已经多余：如果去掉某个未去的站，已扫描的站和其余未去的站仍能通过覆盖核验，就把它剪掉；每次只剪一个，剪完再核验。未去的站也可以挪：把它往当前路线的投影点、路线中点或附近已知源的方向试探，只要挪过去以后覆盖核验仍然通过，而且局部路程至少缩短 $1$ m，就采用新位置，见图\ref{fig:q3-cover}(b)。剪站和挪站都只改计划，不改已经扫过的记录，所以“未发现频道判为无源”的依据始终是真实扫描过的位置。

\begin{center}
\begin{minipage}{\linewidth}
\centering
\QThreeCoverFigure
\captionof{figure}{发现覆盖的布局与调整。(a) 半径 $999$ m 圆周上的 $7$ 个检测站，每站的 $1000$ m 接收圆合起来盖住整个场地。(b) 扫过几站之后，负责区域已被其他站盖住的站剪掉，未去的站往路线附近挪，两种调整都要重新通过覆盖核验。}
\label{fig:q3-cover}
\end{minipage}
\end{center}

\textbf{（3）已发现源的定位：选点规则。}某个频道收到第一次示向度后，定位区域是一条长约 $1500$ m 的窄带，处理它时按第二问的思路继续选点测向。本问对第二问的候选点作了调整，以适应多次测向和从任意位置出发的情形。设当前位置为 $x_0$，区域包围圆的圆心为 $z$、半径为 $r$，$u$ 为从 $x_0$ 指向 $z$ 的单位向量，$v$ 与 $u$ 垂直。候选点取圆心 $z$、当前位置 $x_0$，以及
\[
z+r\,(a\,u+b\,v),\qquad a\in\{0,\pm0.35,\pm0.7\},\ b\in\{0,\pm0.15,\pm0.4\}
\]
中除圆心以外的 $24$ 个点，共至多 $26$ 个，见图\ref{fig:q3-probe}(a)。横向偏移 $b$ 的作用和第二问一样，是拉开交会角；纵向偏移 $a$ 让候选点有靠近当前位置、少走路的选择。去掉该频道已经测过的位置，再去掉不能保证收到信号的点，即到区域任一顶点的距离达到 $999.9$ m 的点。

每个候选点 $x$ 的代价仍是“移动时间加测后剩余代价”。测后区域取决于源的真实位置和这次的角误差，两者都未知，处理办法是取平均：源的位置按面积在区域内取样（把多边形剖成三角形，每个三角形取三个内点，按面积加权），角误差取 $-0.77^\circ$、$0$、$+0.77^\circ$ 三个值，按 $5/18$、$4/9$、$5/18$ 加权，这是 $\pm1^\circ$ 区间上三点求积的节点和权重。对每一对假想的源位置 $p$ 和角误差，把区域与新楔形求交，得到新的包围圆圆心 $z'$ 和半径 $r'$，剩余代价按下式估计：
\begin{equation}\label{eq:q3-4}
P(x)=
\begin{cases}
\dfrac{\max\{0,\ \|z'-x\|-(20-r')\}}{v}+5, & r'\le19.5\ \text{m},\\[8pt]
\dfrac{\|z'-x\|+\max\{0,\ \|p-z'\|-19.5\}}{v}
+6\log_2\dfrac{r'}{19.5}+\dfrac{0.5\,r'}{v}+5, & r'>19.5\ \text{m}.
\end{cases}
\end{equation}
上一行是测完就能清除的情形：走到清除位置再清除。下一行是还要继续测的情形：走到新区域中心、再走到源的时间，加上后续测向的时间，按“每测一次半径大约减半、每次检测加换频 $6$ s”估计为 $6\log_2(r'/19.5)$，再加一项与半径成比例的额外移动。代价取所有假想情形的加权平均，再加上从 $x_0$ 到 $x$ 的移动时间和 $5$ s 检测，选代价最小的候选点。

这一套选点每个源最多做 $4$ 次；超过 $4$ 次，或者包围圆半径已经小于 $22$ m，就直接走到包围圆圆心测一次。在圆心测向有确定的效果：源在圆心的某一侧，楔形从圆心出发，新区域落在一条从圆心出发、长不超过 $r$ 的窄条里，包围圆半径至少减半。所以无论前面的候选点选得好不好，一个源总能在有限次测向内达到清除条件。圆心离区域里每个点都不超过 $r$，远小于 $1000$ m，在那里一定能收到信号。

\textbf{（4）清除判据与清除位置。}光学清除要求离源不超过 $20$ m。定位区域的包围圆半径 $r\le r_c=19.5$ m 时，所有到区域每个顶点都不超过 $20$ m 的位置构成一个非空的凸集，圆心就在其中；走到其中任何一点清除都一定成功。程序在这个集合里选一个点，使“从当前位置走过去”加上“从那里走到下一个任务”的总路程最短：候选包括从当前位置和从下一任务位置各自最近的可行点、两者连线上的插值点，以及当前位置到下一任务的线段与可行集的交（若相交，清除就不需要绕路）。此外，每一步开始时先检查当前位置：若它离某个已知源区域的所有顶点都不超过 $20$ m，就地清除，不再移动。

\textbf{（5）提前光学清除。}上面的清除判据是充分条件，要求区域缩到半径 $19.5$ m 以内，这往往需要在很近的距离上再测一两次。一次失败的光学尝试只花 $3$ s，比一次检测加换频的 $6$ s 还便宜，而且机器狗常常已经站在区域附近。于是本文加了一条规则：区域还没有小到能保证清除，但一个 $20$ m 圆盘已经能盖住它的大部分时，先试一次光学清除，失败了再按原来的流程测。

“盖住大部分”用面积占比来衡量。把定位区域剖成三角形，每个三角形再分成 $10\times10$ 个等面积的小三角形，取各自的重心作为面积样本点 $x_j$，权 $w_j$ 与所在三角形的面积成正比、总和为 $1$。位置 $q$ 处 $20$ m 圆盘盖住的面积占比为
\begin{equation}\label{eq:q3-14}
h_f(q)=\sum_j w_j\,\mathbf 1\{\|x_j-q\|\le20\},
\end{equation}
见图\ref{fig:q3-probe}(b)。它只是区域面积的比例，不是源落在圆盘内的概率，因为源在区域内的分布并不知道；用它来决定要不要花 $3$ s 赌一次，是够用的。提前尝试的条件有四条：该频道正在定位中，且包围圆半径不超过 $80$ m；对这个频道的提前尝试还不到 $2$ 次，成功和失败都计数；当前坐标没有对它试过；$h_f(\text{当前位置})\ge0.4$。四条都满足才向设备发出清除请求。请求成功，该频道转入已清除；失败，定位区域和所有状态原样保留，不做任何推断，因为失败的原因可能是源在圆盘外的那部分区域里，而这部分区域本来就没有被排除。

提前尝试还改变选点：只要一个源满足前三条资格、局部测向还不到 $4$ 次，就不再按（3）的代价选点，而是走到区域的面积加权重心 $q^*=\sum_j w_jx_j$，条件是 $h_f(q^*)\ge0.4$、那里能保证收到信号、且没有对它试过。到达以后先试一次清除；失败再在原地测向、更新区域，然后继续。每源最多两次尝试，所以失败的光学费用每源不超过 $6$ s；但去重心和去代价最小点走的路不同，总的影响要靠实验看。

\begin{center}
\begin{minipage}{\linewidth}
\centering
\QThreeProbeFigure
\captionof{figure}{定位阶段的两种选点。(a) 在包围圆圆心附近按纵向、横向偏移布设候选探测点，超出可靠接收范围或已测过的点剔除，按式\eqref{eq:q3-4} 的代价取最小。(b) 区域已经较小时，在面积加权重心处一个 $20$ m 圆盘盖住的面积占比 $h_f$ 达到 $0.4$，就先试一次光学清除。}
\label{fig:q3-probe}
\end{minipage}
\end{center}

\textbf{（6）任务的安排与共享测量。}每一步的任务集由两部分组成：还没去的检测站，位置就是站的坐标；已发现但没清除的源，位置用当前包围圆的圆心代表。把当前位置当起点，对这些任务求一条不必返回的最短路线，只执行路线上的第一个任务，并记下第二个任务的位置供（4）选清除点用；执行完一个任务，区域和任务集都变了，再重新排。任务最多 $7+16$ 个，程序用多起点的局部搜索求路线：以每个任务为第一站做最近邻排序，各自用 2-opt 和单点移位改进到不能再改，再对最好的一条随机交换两三对任务后重新改进，重复 $40$ 次，取最短的。每次规划只要几毫秒。

机器狗到一个站扫完未发现频道之后，人已经在那里，对附近的已知源顺便再测一次往往划算。程序对每个正在定位的源做一次估计：假设从当前位置对它测向、示向度正指区域中心，算出新区域的包围圆半径；若能缩到原来的 $0.8$ 倍以下并且至少缩小 $30$ m，或者直接缩到清除门限以内，就补测一次，前提是当前位置离这个源不超过 $1500$ m 加半径，且没有在这里测过。清除一个源以后、对某个源测向以后，也用同样的规则检查其他源。另外，如果在当前位置把未发现频道扫一遍能让某个计划中的站被剪掉，就在这里扫，省去以后专门跑一趟。

\textbf{（7）终止条件。}任务在两种情况下结束：已经清除了 $16$ 个不同频道，源数上限保证没有遗漏；或者计划中的站全部扫完（含被剪掉的站，剪站时已核验其余站足以覆盖），此时每个未发现频道在场地任何位置都会被某个已扫描的站收到，可以判为无源，再把已发现的源全部清除。程序在两种情况下都检查 $20$ 个频道是否全部落在已清除或无源两类，缺一个都不退出。发现的源已达 $16$ 个但还有源没清除，或者站扫完了还有已知源，都不能结束。
%
\subsubsection{模型的求解}\label{sec:q3-3}

模型里没有一个可以一次解出的优化问题：覆盖核验、区域更新和清除判据是确定的几何计算，选探测点和排路线各是一个小规模的搜索，把它们按规则串起来、每执行一步用真实反馈更新一次，就是求解过程。流程见图\ref{fig:q3-flow}，步骤归纳为算法\ref{alg:q3-search}。整个过程只用四种真实反馈：示向度、$5$ m 内的近场反馈、无信号，以及清除成功与否。几何部分保证不漏源、不裁掉真实源、判定能清除时一定能清除；选点代价、路线搜索和提前光学的面积占比只影响快慢，算错了最多多花时间，不会漏掉源。

\begin{center}
\begin{minipage}{\linewidth}
\centering
\QThreeFlowFigure
\captionof{figure}{问题三的求解流程。每执行一个任务就重新排路线；处理源时先看能否直接清除，再看是否满足提前光学的条件，最后才选点测向。}
\label{fig:q3-flow}
\end{minipage}
\end{center}

\par\medskip
\noindent\begin{minipage}{\linewidth}
\small
\refstepcounter{algorithm}\label{alg:q3-search}
\textbf{算法\thealgorithm：多源搜索与清除的滚动规划}\par\smallskip
\noindent\textbf{输入：}目标圆域 $\Omega$、$20$ 个频道、源数范围 $10$ 到 $16$、$7$ 个检测站、设备动作费用
\begin{enumerate}
\setlength{\itemsep}{2pt}
\setlength{\parskip}{0pt}
\item 从原点出发，$20$ 个频道全部记为未发现，$7$ 个站全部记为未去；
\item\label{step:q3-check} 对每个已知源，若当前位置离其区域所有顶点都不超过 $20$ m，就地清除；若已清除 $16$ 个频道，或未去的站为空且已知源全部清除，结束；
\item 若还有未发现频道，尝试剪掉多余的站、把未去的站往路线上挪，每次改动都重新通过覆盖核验；
\item 把未去的站和已知源（用包围圆圆心代表）排成从当前位置出发的最短路线，取第一个任务，记下第二个任务的位置；
\item 首任务是检测站：前往该站，从当前频道开始把未发现频道逐个测一遍；对满足补测条件的已知源再测一次；把该站记为已扫描；若未发现频道在所有已扫描站都无信号且未去的站为空，判这些频道无源；回到步骤\ref{step:q3-check}；
\item 首任务是源：若包围圆半径不超过 $19.5$ m，走到兼顾下一任务的清除点清除；否则，若该源满足提前光学的资格，走到区域的面积加权重心先试一次清除，失败则原地测向；否则按式\eqref{eq:q3-4} 在候选点中选代价最小的点测向（局部测向已满 $4$ 次或半径小于 $22$ m 时改测圆心）；用真实示向度和无信号记录更新区域；更新后若半径不超过 $19.5$ m 则立即清除；对满足补测条件的其他已知源再测一次；回到步骤\ref{step:q3-check}。
\end{enumerate}
\end{minipage}
\par\medskip

程序用 Python 实现，几何部分只用 NumPy，选点代价和路线搜索用 Numba 编译。在本地模拟器上，一局从进入到退出的计算时间平均 $0.057$ s，远小于机器狗的动作时间，因此每一步都重新规划不构成负担。运行中有三处检查：在圆心测向却收不到信号、更新后的区域为空、判定能清除却清除失败，任何一处出现都说明前面的推理和实际反馈矛盾，程序立即报错停机，不会把矛盾的状态继续用下去；主批 $10000$ 次运行中这三处检查一次也没有触发。
%
\subsubsection{结果分析}\label{sec:q3-4}

\textbf{（1）实验设置。}策略的效果在本地复原的模拟器上检验。这个模拟器按题目附件 2 的接口说明和演练程序的生成规则复原：源在半径 $1770$ m 的圆内按面积均匀分布，数量在 $10$ 到 $16$ 之间等概率，频道从 $20$ 个中不放回抽取，接收半径在 $1000$ 到 $1500$ m 之间均匀抽取；测向误差是在 $150$ m 网格上平滑变化的空间误差场，先量化到 $0.01^\circ$ 再限制在 $\pm1^\circ$ 内，同一位置、同一频道的误差固定，重复检测不能消除；计时按题目规则，每段移动的时间取整到微秒。它不是官方程序本身，没有和官方程序做过逐输入输出的对照，因此这里的数字不能当作正式测试成绩，只用来在相同场景上比较不同策略。

比较五种策略，都由同一套几何模块搭成，区别只在规则：
\begin{itemize}
\item \textbf{本文方案}：即算法\ref{alg:q3-search}，含提前光学清除；
\item \textbf{无提前光学}：去掉（5）的规则，区域必须缩到 $19.5$ m 以内才清除，其余相同；
\item \textbf{原点先扫}：在无提前光学的基础上，开局先在原点把 $20$ 个频道扫一遍；
\item \textbf{联合规划}：在无提前光学的基础上，排路线时把已知源的位置也当作可能的扫描点，若能替代某个站就剪掉那个站；
\item \textbf{早期版本}：候选探测点只在圆心两侧横向取 $6$ 个、路线只用最近邻加 2-opt 的更早一版。
\end{itemize}
场景共 $2000$ 个，每个场景五种策略从同一初始状态出发，源位置、接收半径和误差场完全相同，共 $10000$ 次运行；场景在策略冻结之后才生成，没有用来调参数。评价指标是每个场景的总耗时除以该场景的源数，再对场景等权平均；总耗时包括清除最后一个源之后为确认没有遗漏而做的扫描。策略之间的差用同场景配对的 $10000$ 次自助法给出 $95\%$ 区间。另用 $50$ 个场景、每种策略各跑两遍、单进程顺序执行，测每局的计算时间。这五种策略是从更多候选中按此前 $184$ 局官方演练的成绩挑出来的，但那些演练没有让各策略跑相同的场景，不能用来排名，所以本文的比较以下面 $2000$ 个同场景的结果为准。

\textbf{（2）主要结果。}$10000$ 次运行全部完整清除，共 $26006$ 个源实例，没有失败或超时的场景。各策略的平均定位清除时间见表\ref{tab:q3-main}，相对无提前光学方案的同场景差值见表\ref{tab:q3-paired}。

\Needspace{46mm}
\begin{center}
\begin{minipage}{\linewidth}
\centering\small
\captionof{table}{五种策略在 $2000$ 个相同场景上的平均定位清除时间（单位：s/源）}
\label{tab:q3-main}
\vspace{4pt}
\begin{tabular}{@{}lrcrrrr@{}}
\toprule
策略 & 均值 & 均值 95\% 区间 & 中位数 & P95 & 全清场景 & 计算时间 / s/局 \\
\midrule
本文方案 & $239.26$ & {[}237.85, 240.68{]} & $236.60$ & $296.29$ & $2000/2000$ & $0.057$ \\
无提前光学 & $240.89$ & {[}239.45, 242.36{]} & $238.07$ & $297.68$ & $2000/2000$ & $0.049$ \\
联合规划 & $240.92$ & {[}239.46, 242.38{]} & $238.17$ & $297.76$ & $2000/2000$ & $0.066$ \\
原点先扫 & $241.38$ & {[}239.87, 242.87{]} & $238.81$ & $300.39$ & $2000/2000$ & $0.056$ \\
早期版本 & $243.85$ & {[}242.34, 245.36{]} & $241.62$ & $303.70$ & $2000/2000$ & $0.125$ \\
\bottomrule
\end{tabular}
\end{minipage}
\end{center}

\Needspace{40mm}
\begin{center}
\begin{minipage}{\linewidth}
\centering\small
\captionof{table}{相对无提前光学方案的同场景差值（负值表示更快，单位：s/源）}
\label{tab:q3-paired}
\vspace{4pt}
\begin{tabular}{@{}lrcr@{}}
\toprule
策略 & 平均差值 & 差值 95\% 区间 & 更快 / 相同 / 更慢的场景数 \\
\midrule
本文方案 & $-1.631$ & {[}$-1.715$, $-1.546${]} & $1736\ /\ 0\ /\ 264$ \\
联合规划 & $+0.032$ & {[}$-0.019$, $+0.084${]} & $86\ /\ 1821\ /\ 93$ \\
原点先扫 & $+0.493$ & {[}$+0.034$, $+0.946${]} & $955\ /\ 0\ /\ 1045$ \\
早期版本 & $+2.958$ & {[}$+2.477$, $+3.433${]} & $760\ /\ 0\ /\ 1240$ \\
\bottomrule
\end{tabular}
\end{minipage}
\end{center}

本文方案的均值最低，比无提前光学方案每源平均省 $1.63$ s，区间不含零；$2000$ 个场景中 $1736$ 个更快、$264$ 个更慢，没有一个相同，说明提前光学几乎在每一局都改变了动作序列。另外三种对照各自说明一件事：联合规划和无提前光学在 $1821$ 个场景上动作完全相同，其余场景互有快慢，平均差值的区间跨零，也就是把已知源位置当扫描点来剪站，在这个问题上基本没有机会用上；原点先扫平均慢 $0.49$ s/源，区间刚好不含零，胜负场景接近一半对一半，说明开局那 $119$ s 的检测费用大体上被早发现源省下的路程抵消，但抵消得不够；早期版本平均慢 $2.96$ s/源，是候选探测点和路线搜索两处改进的合计收益。

\textbf{（3）时间省在哪里。}把本文方案和无提前光学方案每局的动作分开看：平均检测次数由 $123.8$ 次降到 $118.0$ 次，少了 $5.7$ 次，约 $29$ s；平均移动距离由 $11358$ m 增到 $11410$ m，多走 $52$ m，约 $10$ s；提前尝试在 $2000$ 局里共失败 $1616$ 次，平均每局 $0.8$ 次，约 $2.4$ s。三项合计与每局约 $21$ s 的实际节省相符（每局平均 $13.0$ 个源）。也就是说，提前光学省下的是本来要在近距离上再做的一两次测向；代价是失败的尝试和去重心多走的路，两者都比省下的检测便宜。这 $1616$ 次失败全部按 $3$ s 计入总耗时，没有从结果里去掉。

\textbf{（4）按源数分组。}表\ref{tab:q3-count} 和图\ref{fig:q3-strata} 按场景的真实源数（只在事后分组时使用）列出结果。源越多，每源平均时间越短，因为扫站的固定开销被更多的源分摊；源数为 $16$ 的场景还能在发现第 $16$ 个源后停止扫站。本文方案在七个分组里都比无提前光学快，差值稳定在 $1.5$ 到 $1.8$ s/源之间，与源数无关；原点先扫在源少的分组里更慢、源多的分组里略快，因为源多时开局多扫一遍更容易早发现几个源。

\Needspace{44mm}
\begin{center}
\begin{minipage}{\linewidth}
\centering\small
\captionof{table}{按源数分组的平均定位清除时间（单位：s/源）}
\label{tab:q3-count}
\vspace{4pt}
\begin{tabular}{@{}lrrrrrrr@{}}
\toprule
源数 & 10 & 11 & 12 & 13 & 14 & 15 & 16 \\
\midrule
场景数 & 295 & 268 & 299 & 291 & 259 & 297 & 291 \\
本文方案 & $286.95$ & $267.45$ & $250.87$ & $235.47$ & $223.15$ & $213.51$ & $197.43$ \\
无提前光学 & $288.48$ & $269.11$ & $252.52$ & $237.03$ & $224.78$ & $215.14$ & $199.20$ \\
联合规划 & $288.52$ & $269.28$ & $252.56$ & $237.07$ & $224.70$ & $215.19$ & $199.17$ \\
原点先扫 & $290.46$ & $270.70$ & $252.73$ & $238.64$ & $224.48$ & $214.17$ & $198.54$ \\
早期版本 & $293.16$ & $273.16$ & $255.03$ & $241.29$ & $226.90$ & $216.67$ & $200.77$ \\
\bottomrule
\end{tabular}
\end{minipage}
\end{center}

\begin{center}
\begin{minipage}{\linewidth}
\centering
\QThreeStrataFigure
\captionof{figure}{按源数分组的结果。(a) 四种策略的平均定位清除时间都随源数增加而下降。(b) 相对无提前光学方案的同场景差值，本文方案在每个分组都在零线以下，联合规划贴着零线。}
\label{fig:q3-strata}
\end{minipage}
\end{center}

\textbf{（5）个例。}平均改善不等于每个场景都改善。以逐例记录中的几个场景为例：编号 f392164b 的 $11$ 源场景，本文方案 $239.83$ s/源，无提前光学 $249.36$ s/源，省了 $9.5$ s/源，是提前尝试连续命中的情形；编号 2baf85fc 的 $11$ 源场景，本文方案 $274.28$ s/源，比无提前光学的 $274.10$ 慢 $0.18$ s/源，是一两次尝试落空、又多走了一段路的情形。慢的 $264$ 个场景里，差值多数在 $1$ s/源以内，这与每源最多两次、每次 $3$ s 的费用上限一致；去重心测向引起的路线变化才是更慢场景的主要来源，这一部分没有上限，但从分布看很少超过几秒。

\textbf{（6）计算开销。}表\ref{tab:q3-main} 最后一列是单进程顺序复跑 $50$ 个场景得到的每局计算时间，本文方案 $0.057$ s，比无提前光学多 $0.008$ s，多出的部分是面积样本点的生成和占比计算；早期版本最慢，因为它的路线搜索没有编译。计算时间都远小于机器狗一局几千秒的动作时间。$10000$ 次运行的每一个动作都按记录重放核对过，反馈和逐段计时一致；每个场景都检查了退出条件，源数少于 $16$ 时所有未发现频道在每个已扫描站都无信号且未去的站为空，源数为 $16$ 时有 $16$ 次真实的清除成功。

\textbf{（7）适用条件。}以上结果只对本文冻结的参数和本地模拟器的场景生成方式成立。提前光学的门槛 $0.4$、$80$ m 和 $2$ 次是开发阶段定下的，没有按这 $2000$ 个场景重新调过；如果官方演练的误差特性与此差别较大，提前尝试的命中率可能变化，但它的费用上限不变，几何部分的保证也不受影响：无论尝试成功与否，站扫完就不会漏源，判定能清除时一定能清除。官方演练和正式测试的清除比例与平均定位清除时间，要以实际日志为准。
%
\subsubsection{问题三的结论}\label{sec:q3-5}

把本节的做法归纳成对题目的回答：

\begin{itemize}
\item \textbf{搜索策略。}用半径 $999$ m 圆周上等角分布的 $7$ 个检测站保证发现：每站的 $1000$ m 接收圆合起来盖住整个场地，任何位置的全向源都会在某个站被收到，因此不必在原点先扫全部频道。每到一站把未发现频道扫一遍，顺便对附近的已知源补测；扫过几站后，多余的站剪掉、未去的站往路线上挪，每次调整都重新核验覆盖。
\item \textbf{定位策略。}每个已发现的源维护一个凸多边形定位区域，示向度交楔形，无信号点挖去 $1000$ m 圆盘并与有信号点配对切半平面。继续测向时在包围圆圆心附近布设候选点，按移动时间加测后剩余代价选点，最多 $4$ 次后回圆心测，每次至少把半径减半。
\item \textbf{清除策略。}包围圆半径不超过 $19.5$ m 时，走到离区域所有顶点都不超过 $20$ m、且兼顾下一任务的位置清除，一定成功。半径不超过 $80$ m、$20$ m 圆盘能盖住区域面积四成以上时，先试一次光学清除，每源最多两次，失败只花 $3$ s 且不改变区域。只有收到真实的成功反馈才算清除。
\item \textbf{任务安排。}未去的站和已知源合成一个路线问题，用多起点局部搜索排序，只执行第一个任务后重新规划；$16$ 个频道全部清除，或站扫完且已知源全部清除时结束。
\end{itemize}

在本地复原模拟器的 $2000$ 个场景上，这套策略全部完成清除，平均定位清除时间 $239.3$ s/源，比不做提前光学清除的方案省 $1.63$ s/源，在 $1736$ 个场景上更快。它省时间的来源是用便宜的光学尝试替代近距离的重复测向；不漏源、不裁掉真实源、判定能清除时一定能清除这三条，不依赖任何尝试的成败。
%
